\chapter{Лекция}
\label{ch:intro}

\section{Символьная синхронизация}

Символьная синхронизация - процесс поиска (определения) оптимального времени взятия отсчетов принятого сигнала. Она определяет
задержку (сдвиг отсчетов сигнала) перед процедурой downsampling, которая позволяет "попасть" в оптимальные моменты взятия отсчетов.


\section{Архитектура приемника}

Доработаем схему приемника, которую мы сделали на прошлом занятии, добавив блоки, выполняющие символьную синхронизацию.

\begin{figure}[H]
    \centering
    \includegraphics[width=1.0\textwidth]{arch.png}
    \caption{Архитектура приемника}
\end{figure}

\section{Принцип работы символьной синхронизации}
Между согласованным фильтром и процедурой downsampling добавили контур, в котором будет выполняться символьная синхронизация. \\

\section{Timing correction}

Семплы, выходящие из согласованного фильтра поступают на блок Timing corrector, где скорректируются таким образом, чтобы правильно 
начать выбирать символы. 

\section{TED}

Принцип работы символьной синхронизации основан на формировании сигнала ошибки $e(k)$ при принятым отсчетам и коррекции временной
задержки. \\

Задача заключается в том, чтобы найти моменты времени, когда производная от сигналов равна 0 (моменты смены знака). Будем использовать
метод Гарднера. Он основан на разнице между отсчетами с интервалом в 1 символ. Рассчитать $e(k)$ по схеме Гарднера можно следующим образом:

$$e(k) = (s[k] - s[k-2]) * s[k-1]$$

Здесь $s[k]$ - определенный отсчет сигнала.\\

Визуализация метода:

\begin{figure}[H]
    \centering
    \includegraphics[width=1.0\textwidth]{scheme.png}
    \caption{Расчет TED}
\end{figure}

Если визуализировать коэффициенты TED, то должна получиться кривая, схожая с sin или cos

\begin{figure}[H]
    \centering
    \includegraphics[width=1.0\textwidth]{e.png}
    \caption{График TED}
\end{figure}

\section{Loop filter}

Этот блок по определенному алгоритму усредняет (сглаживает) TED и вычисляет сдвиг. \\

Схема работы Loop filter

\begin{figure}[H]
    \centering
    \includegraphics[width=1.0\textwidth]{scheme2.png}
    \caption{Схема Loop filter}
\end{figure}

Формулы для вычисления K1, K2

\begin{figure}[H]
    \centering
    \includegraphics[width=1.0\textwidth]{forms.png}
    \caption{Формулы}
\end{figure}

В вычислениях будем использовать следующие константы (получены экспериментальным путем)

\begin{lstlisting}
    BnTs = 0.01
    Nsps = 10
    Kp = 0.002
    zeta = sqrt(2) / 2
\end{lstlisting}

\endinput
